\begin{scanpage}{0013}{ix}
Let a plane electromagnetic wave fall on some ideally conducting body which is found in free space. In the physical optics approach, the surface current density which is induced by this wave on the irradiated part of the body's surface is (in the absolute system of units) equal to
\[\mathbf j^0=\frac c{2\pi}[\mathbf n\mathbf H_0].\tag{A}\]
where $c$ is the speed of light in a vacuum, $\mathbf n$ is the external normal to the body's surface, $\mathbf H_0$ is the magnetic field of the incident wave. On the darkened side of the body the surface current is assumed to be equal to zero $(\mathbf j^0=0)$. Equation (A) means that on each element of the body's irradiated surface the same current is excited as on an ideally conducting surface of infinite dimensions tangent to this element. The scattered field created by the current (A) is then found by means of Maxwell's equations.

It is obvious that in reality the current induced on the body's surface will differ (as a consequence of the curve of the surface) from the current $\mathbf j^0$. The precise expression for the surface current density has the form
\[\mathbf j=\mathbf j^0+\mathbf j^1.\tag{B}\]
where $\mathbf j^1$ is the surface density of the additional current which results from the curve of the surface. By the curve of the surface, we mean any of its deviations from an infinite plane (a smooth curve, a sharp bend, a bulge, a hole, etc.). If the body is convex and smooth and its dimensions and radii of curvature are large in comparison with the wavelength, then the additional current is concentrated mainly in the vicinity of the boundary between the illuminated and shadowed parts of the body's surface. But if the body has an edge, bend, or point, then the additional current also arises near them. The additional current density is comparable to the density $\mathbf j^0$, as a rule, only at distances of the order of a wavelength from the corresponding edge, bend, or point. Thus, if the body's dimensions
\end{scanpage}
